\documentclass[a4paper, 11pt]{article}
\usepackage[utf8]{inputenc}
\usepackage[T1]{fontenc}
\usepackage{polski}


\usepackage{times}
\usepackage[margin=2.5cm]{geometry}

\usepackage{apacite}
\usepackage{natbib}

\author{Imię Nazwisko\\Nr indeksu\\E-mail}
\title{Tytuł pracy}
\date{\today}

\begin{document}
\maketitle

\begin{abstract}
  Tutaj abstrakt (250 słów) ...
  
  
oindent
  \textbf{Słowa kluczowe:} pierwsze, drugie, trzecie, czwarte, piąte
\end{abstract}


ewpage %%%%%%%%%%%%%%%%%%%%%%%%%%%%%%%%%%%%%%%%%%%%%%%%%%%%%%%%%%%%%%

Praca. Cztery strony tekstu. Przykład cytowania, por. \citep{przyklad} (spis literatury w~pliku \textit{bibliografia.bib}). Spis literatury powinien obejmować tylko pozycje wykorzystane w tekście.

Przykłady cytowań:
\begin{itemize}
\item \cite{typografia}
\item \citep{typografia}
\item \cite[s.~12]{typografia}
\item \citeauthor{typografia}
\item \citeyear{typografia}
\end{itemize}



ewpage  

%%%%%%%%%%%%%%%%%%%%%%%%%%%
%%%%%%%%%%%%%%%%%%%%%%%%%%%
\bibliographystyle{apacite}
\bibliography{bibliografia}
\end{document}